\chapter{ทบทวนวรรณกรรม}

\section{การตรวจจับความผิดปกติในอนุกรมเวลา}

การตรวจจับความผิดปกติเป็นปัญหาที่เกี่ยวข้องกับการระบุรูปแบบข้อมูลที่เบี่ยงเบนจากพฤติกรรมปกติ \citep{chandola2009anomaly,pang2021deep} สำหรับข้อมูลอนุกรมเวลา ความผิดปกติอาจเกิดขึ้นในรูปแบบจุดเดี่ยว ช่วงเวลา หรือการเปลี่ยนแปลงของ regime

\section{Transformer สำหรับข้อมูลอนุกรมเวลา}

Transformer ใช้กลไก attention เพื่อเรียนรู้ความสัมพันธ์ระหว่างตำแหน่งในลำดับข้อมูล \citep{vaswani2017attention} ในบริบทของ anomaly detection งาน Anomaly Transformer เสนอแนวคิด association discrepancy เพื่อวัดความแตกต่างระหว่างความสัมพันธ์เชิง prior และ series association \citep{xu2022anomaly}

\section{Geometric Brownian Motion}

Geometric Brownian Motion เป็นแบบจำลองพื้นฐานที่ใช้ในคณิตศาสตร์การเงินเพื่ออธิบายพลวัตของราคาสินทรัพย์ภายใต้สมมติฐานผลตอบแทนต่อเนื่องและความผันผวน \citep{merton1973theory,hull2018options} แม้ GBM จะเป็นข้อสมมติที่เรียบง่าย แต่งานวิจัยนี้ใช้ GBM เป็นกรอบอ้างอิงเชิงการแจกแจง ไม่ใช่การอ้างว่าตลาดหุ้นเป็นไปตาม GBM อย่างสมบูรณ์

\section{ช่องว่างของงานวิจัย}

ช่องว่างสำคัญคือการเชื่อมแบบจำลอง deep learning สำหรับอนุกรมเวลากับข้อจำกัดเชิงการเงินที่ตีความได้ งานวิจัยนี้จึงศึกษาการใช้ GBM-implied return distribution เป็นองค์ประกอบหนึ่งของ anomaly score เพื่อให้การตรวจจับความผิดปกติเชื่อมโยงกับพฤติกรรมผลตอบแทนและความผันผวนโดยตรงมากขึ้น

